\section{Examples}

Each example below was run with the version of \tool this manual ships with.

\subsection{Reachability and safety in conflict}

\begin{lstlisting}
$ ./build/tctl_sat 'EF[1,2] q & AG[0,3] !q'
1: UNSATISFIABLE
\end{lstlisting}

\noindent
The first conjunct demands $q$ somewhere in $[1,2]$; the second forbids it
everywhere in $[0,3]$.  Since $[1,2]\subseteq[0,3]$ there is no model.  Widen
the safety window so the two no longer overlap and a model appears:

\begin{lstlisting}
$ ./build/tctl_sat 'EF[2,2] q & AG[0,1] !q'
1: SATISFIABLE
\end{lstlisting}

\subsection{Strictness matters}

\begin{lstlisting}
$ printf 'EF[1,2] p & AG[0,1] !p\nEF(1,2) p & AG[0,1] !p\n' > s.txt
$ ./build/tctl_sat s.txt
1: SATISFIABLE
2: SATISFIABLE
\end{lstlisting}

\noindent
Both hold: $p$ can be placed strictly after $1$.  The zone layer tracks the
open/closed endpoint exactly rather than rounding it to an integer.

\subsection{A leads-to specification}

The classic response property, \texttt{AG(request -> AF[0,5] grant)}, is the
shape most often written for \textsc{Uppaal}.  It is satisfiable on its own, and
stays so together with a bounded-response requirement:

\begin{lstlisting}
$ ./build/tctl_sat --tot leadsto.txt
1: SATISFIABLE
\end{lstlisting}

\noindent
where \texttt{leadsto.txt} contains

\begin{lstlisting}[style=formula]
AG (request -> AF[0,5] grant)
AG[0,20] !error
EF[1,3] request
\end{lstlisting}

\subsection{Nested timed modalities}

Nesting is where the procedure does real work.  The following is satisfiable,
and the node count shows the cost:

\begin{lstlisting}
$ ./build/tctl_sat --stats 'EG[1,5] (EF[2,4] (EG[6,8] p))'
1: SATISFIABLE
  Stats: nodes=261
\end{lstlisting}

\subsection{A contradiction the tool must catch}

\begin{lstlisting}
$ ./build/tctl_sat 'EF p & !EF p'
1: UNSATISFIABLE
\end{lstlisting}
